\documentclass[11pt]{amsart}
\usepackage{amsmath,amssymb,amsthm,mathtools}
\usepackage[hidelinks]{hyperref}
\hypersetup{pdftitle={The density conjecture for activated random walk on the square lattice},pdfauthor={Kaiwen Jin}}
\newtheorem{theorem}{Theorem}
\newtheorem{lemma}[theorem]{Lemma}
\newtheorem{proposition}[theorem]{Proposition}
\newcommand{\Z}{\mathbb Z}
\newcommand{\E}{\mathbb E}
\newcommand{\Prob}{\mathbb P}
\newcommand{\1}{\mathbf 1}
\newcommand{\calE}{\mathcal E}
\newcommand{\calG}{\mathcal G}
\newcommand{\TV}{\operatorname{TV}}
\newcommand{\floor}[1]{\left\lfloor #1\right\rfloor}
\newcommand{\ceil}[1]{\left\lceil #1\right\rceil}
\newcommand{\sr}{\zeta_c}
\author{Kaiwen Jin}
\email{yui.ui.siki@gmail.com}
\subjclass[2020]{60K35, 82C22}
\keywords{Activated random walk, stationary density, fixation}
\title[Stationary density of activated random walk]{The density conjecture for activated random walk on the square lattice}
\begin{document}
\begin{abstract}
We prove the density conjecture for activated random walk on the square lattice. For every positive sleep rate, the expected density in the stationary state of a driven killed square converges to the fixation threshold for iid Poisson initial data on the infinite lattice. Comparing the area-normalized logarithms of the sleeper count's Laplace transform across square sizes gives positive probability of density within any fixed tolerance of the threshold. A Dirichlet estimate for the limiting occupation potentials then permits approximation of the influxes on most tiles by a small finite family of profiles. Independent stabilization tests for each member of this family rule out a stationary mean above the threshold.
\end{abstract}
\maketitle
\section{Introduction}

Activated random walk allows moving particles to become immobile and to be reactivated by other particles. Consider the process on a box with absorbing boundary. Repeated additions of one particle, each followed by relaxation, produce a stationary distribution of sleepers. The density conjecture asks whether the density selected by this driven system approaches the critical density of the infinite system started from independent particle counts. The equality was proved on the line by Hoffman, Johnson and Junge \cite{HJJ}. Forien gave another proof using superadditivity \cite{Forien}. We consider the square lattice.

\begin{theorem}\label{t:main}
 Fix a sleep rate \(\lambda>0\). On \(\Z^2\) let active particles jump at rate one to a uniform nearest neighbor. A lone active sleeps at rate \(\lambda\) and an arriving active wakes a sleeper. Let \(\zeta_c(\lambda)\) be the infimum of densities for which the iid Poisson all-active system almost surely fails to fixate, where fixation means finitely many jumps at every site.

 For the square \(\Lambda_n=[-n,n]^2\cap\Z^2\), kill particles on exit and repeatedly add one active particle at an independent uniform site, stabilizing after each addition. If \(\pi_n\) is the stationary law and \(N\) counts its sleepers, then
 \[
 \lim_{n\to\infty}\frac{\E_{\pi_n} N}{|\Lambda_n|}=\zeta_c(\lambda).
 \]
\end{theorem}

Two facts about finite boxes give a starting point. The stationary law is sampled by stabilizing one active particle per site \cite{LL}, and subcritical Poisson input loses only a vanishing fraction through the boundary by the exit criterion of Rolla and Tournier \cite{RT}. They give the lower bound. For the upper bound, compare the actual stabilization in a large square with stabilizations in smaller tiles. The incoming particle counts of a tile depend on the very instructions used inside it. The comparison must handle many possible deterministic arrival profiles before the realized arrivals are selected.

We compare the negative-parameter log moment of the stationary sleeper count per site across square sizes. For each fixed \(\varepsilon>0\) this gives positive probability of a density below \(\zeta_c+\varepsilon\). Simultaneous lower bounds hold for the occupations of smaller squares with probability tending to one. On those two events the rescaled jump field has subsequential limits given by Dirichlet Green potentials. These limits differ from the zero-boundary potential of the constant source \(1-\zeta_c\) by functions with small Hessian in \(L^2\). This gives a finite family of inlet profiles that approximates the influx on most tiles. At a fixed deterministic profile the tile test counts are independent. If the stationary mean on the smaller squares were too large, a union over the profiles would make the low event too rare.

\section{Finite stabilization and the lower density bound}\label{s:finite}

We take a fixed \(0<\lambda<\infty\) throughout. Let \(V_L=[1,L]^2\cap\Z^2\), with \(L\) odd. It is a translate of the square in Theorem~\ref{t:main}. Constants may depend on \(\lambda\), and any further fixed parameters are specified when they are chosen. We write \(\pi_L\) for the stationary configuration law on this square and \(N\) for its number of sleepers.

At each site use independent instructions, each a sleep attempt with probability \(\lambda/(1+\lambda)\) and otherwise a uniform nearest-neighbor jump. Use the next instruction only at an active site. A sleep attempt with two or more actives has no effect. An arrival wakes a sleeper. For graph notation let \(P\) be the discrete jump kernel with all outward jumps sent to one absorbing state. Write \(Q\) for the restriction of \(P\) to the nonsink vertices and \(G=(I-Q)^{-1}\) for the visit Green operator, including time zero. The least-action principle and the abelian property identify the instruction odometers and outputs of complete legal histories on the same stacks \cite[\S1.3]{LL}. In particular, they identify the actual jump counts because each complete history uses the same stack prefixes. The stack law agrees with continuous-time stabilization. One may give each active label additional, ineffective sleep attempts at rate \(\lambda\) while it shares a site, so all active clocks choose the instruction with the same probabilities.

Every nonanticipating legal order until stability is finite almost surely. If \(t(x)\) is the mean exit time of the jump chain, with \(t=0\) at the sink, the sum of \(t\) over particles has mean drift \(-1/(1+\lambda)\) per instruction before stopping. This is the first-step equation \(t-Qt=1\), since a sleep attempt preserves particle locations. Nonnegativity at a deterministic instruction cutoff bounds the expected instruction count by \(1+\lambda\) times the initial sum. The same argument works when some sites are frozen. At such a finite stop, the unused tails are independent with the original stack law \cite[Lemma~1]{LL}. This remains true with random starting data independent of the stacks, by conditioning on those data.

One active at every site will be called the saturated input. Its stabilization has law \(\pi_L\). In fact, the finite-graph exact sampler is invariant under each fixed-site addition, and so under the iid uniform drive \cite[Theorem~1]{LL}. The entire drive is independent of the whole instruction-stack array. Each site is chosen infinitely often almost surely and the sink is accessible, so the recurrent class theorem \cite[Lemma~5]{LL} gives the whole stable state space as the single recurrent class at strictly positive finite rates. Successive unused tails give the same homogeneous driven kernel. Thus the invariant law is unique. For later use, the stabilization of saturation plus any fixed deterministic nonnegative vector also has law \(\pi_L\). First move the extra labels to the sink. Each extra label shares its current site with a reserved active, so its sleep attempts are ineffective. Then stabilize the reserved actives on the fresh tails.

We will use the following local form of the abelian property. In a completed ambient history, restrict the sequence of instructions to a square \(K\). Kill every jump leaving \(K\), and treat every arrival from outside \(K\) as an addition. Move all such additions to the beginning. This preserves legality and the terminal state. An addition commutes with an already legal local instruction, including a sleep attempt on a singleton that is then woken by the addition. The resulting local history is complete. By least action, the restriction of the ambient stable state is therefore the killed stabilization on its original \(K\) stacks from the restricted input plus the actual inlet. The inlet vector counts all jumps, including reentries, along each directed incoming boundary edge, called a port. If two ports share a receiving corner, add their counts when forming the local input.

We give the stochastic count comparison on a killed square. Forien proves the corresponding comparison on finite subsets of the line \cite[Lemma~7]{Forien}. For all-active input \(\sigma\), increasing \(f\), and a site \(x\) including the sink, let \(F_f(x)\) be the mean of \(f\) of the retained count from \(\sigma+\delta_x\) minus that from \(\sigma\). Addition at the sink has no effect. Freeze \(x\) and stabilize off \(x\) to get \(\beta_x\). Let \(D_x\) denote the conditional mean of \(f\) after adding to \(\beta_x\) at a fresh destination of law \(P(x,\cdot)\) and stabilizing. An added particle at \(x\) gives this continuation whenever it jumps. If \(\beta_x(x)>0\), wait to the first jump at \(x\), since all preceding sleep attempts there are ineffective. If \(\beta_x(x)=0\), the extra particle can also sleep immediately, leaving \(N(\beta_x)+1\) particles. Averaging the destination term over the off-\(x\) history is equivalent to placing the extra particle at the destination from the start, by commuting the addition with the off-\(x\) prefix. Hence
\[
 (I-Q)F_f(x)=\frac{\lambda}{1+\lambda}\,
     \E\big[\1_{\{\beta_x(x)=0\}}\big(f(N(\beta_x)+1)-D_x\big)\big].
\]
The bracket is evaluated on the displayed event, where it is nonnegative by conservation. The nonnegative inverse \(G\) gives \(F_f\ge0\), the desired stochastic comparison. To compare Poisson input with saturation, first topple only sites with at least two actives, as in \cite[\S6.2]{Forien}. The output of this legal prefix has at most one active per site and all other sites empty. Its continuation on fresh tails is dominated by the saturated count.

For iid Poisson input of density \(0<\rho<\sr(\lambda)\), the expected number of exits per site tends to zero. This is the contrapositive of the finite-box exit criterion \cite[Proposition~3 and \S5]{RT}: a positive limsup of the expected number of exits per site implies almost-sure nonfixation. The criterion covers the symmetric walk and iid initial counts of finite mean. In a finite box, freezing a label outside at its first exit gives the same interior evolution and exit count as killing it. By conservation, Poisson concentration and Markov's inequality, the density retained from this input tends to \(\rho\) in probability and in mean. For the corresponding saturation and exit argument on the line, see \cite[\S5.1]{Forien}. The stochastic comparison and \(\rho\uparrow\sr(\lambda)\) yield, with \(\zeta=\sr(\lambda)\),
\begin{equation}\label{e:lower}
 \liminf_{L\to\infty}L^{-2}\E_{\pi_L}N\ge\zeta,
 \qquad
 \pi_L(N/L^2<\zeta-\varepsilon)\longrightarrow0
 \quad(\varepsilon>0).
\end{equation}
The set of nonfixating Poisson densities is nonempty because at \(\rho>1\), conservation forces an expected number of exits per site at least \(\rho-1\), so \(0\le\zeta\le1\). When \(\zeta=0\), \eqref{e:lower} follows from nonnegativity. Here and below the definition by finite site jump counts agrees with the standard fixation event. To see the continuous-time point, give each label a prospective rate-one walk run in its active time. Over a bounded absolute interval \([0,T]\), the expected total number of prospective visits to any fixed site, using active time up to \(T\), is \(\rho(1+T)\), from the column sums of the full walk kernel and the mean Poisson clock count. Thus actual jumps are locally finite on bounded intervals. One may use independent real-time jump clocks with sleeping ticks discarded to see that jump times do not coincide almost surely. Eventual sitewise constancy then forbids later jumps. Finitely many jumps incident to a site conversely allow only its final sleep. The finite-degree graph lets one use either departure counts or all incident counts in the criterion over all sites.

\section{Pressure and fixation}\label{s:pressure}

Write \(\Phi_L(t)=L^{-2}\log\E_{\pi_L}e^{tN}\). For each fixed real \(t\), the limit \(\Phi(t)\) exists along odd \(L\). Here is the tiling estimate. In the ambient saturated sample, use disjoint odd side-\(\ell\) tiles and leave strips of total area at most \(2L\ell\). The initial exit potential is \(O(L^4)\) by the one-coordinate exit-time bound \(O(L^2)\). Its drift-corrected martingale has increments \(O(L^2)\). At a deterministic cutoff \(H=\ceil{C_{\lambda,t} L^6}\) the event that stabilization is still running forces a martingale increase of at least \(H/[2(1+\lambda)]\) for large \(L\). The exponential bound for bounded martingale increments gives probability at most \(e^{-(2|t|+1)L^2}\) with \(C_{\lambda,t}\) large enough. Let \(T\) be the completed instruction count and \(j=\floor{L/\ell}^2\) the number of tiles. On \(T\le H\) there are at most \((H+1)^{4\ell}\) candidates for each tile inlet. For any deterministic joint row we may drop the event that it is the actual row, then use independent tile counts from \(\pi_\ell\). Thus
\[
 \E_{\pi_L}[e^{tN};\,T\le H]
 \le (H+1)^{4j\ell}e^{|t|(L^2-j\ell^2)}
       \big(\E_{\pi_\ell} e^{tN}\big)^j .
\]
The right side is at least \(e^{-|t|L^2}\), while the weighted \(T>H\) term is at most \(e^{-(|t|+1)L^2}\). Taking logarithms gives
\begin{equation}\label{e:fixed}
 \Phi_L(t)\le\Phi_\ell(t)
   +C_{\lambda,t}\frac{1+\log L}{\ell}
   +2C|t|\frac{\ell}{L}+L^{-2}\log2 .
\end{equation}
The strip term can be dropped for exact multiples. Fix an odd base \(\ell_0\) and iterate on \(\ell_{i+1}=\ell_i^2\). The accumulated error is \(O_{\lambda,t}((1+\log\ell_0)/\ell_0)\). For general large \(L\), choose \(i\) with \(\ell_i^2\le L\le\ell_i^4\) and use that same tile size \(\ell_i\) in \eqref{e:fixed}. The resulting extra error is \(O_{\lambda,t}((1+\log\ell_i)/\ell_i)\). Let \(\ell_0\) run through sizes realizing the liminf. This proves the pressure limit. It is convex, nondecreasing and \(1\)-Lipschitz, with \(\Phi(0)=0\). Set \(u_-=\Phi'(0-)\). At any \(v<u_-\) a fixed \(t<0\) with \(\Phi(t)<tv\) gives a volume-exponential upper bound for \(\pi_\ell(N<v\ell^2)\).

\begin{lemma}\label{l:fix}
 Suppose that for some \(v>0\), \(c>0\) and \(\alpha>0\) the stationary lower tail satisfies \(\pi_\ell(N<v\ell^2)\le e^{-c\ell^{1+\alpha}}\) at all sufficiently large odd \(\ell\). Then iid Poisson input at every \(0<\rho<v\) fixates. In particular \(u_-\le\zeta\).
\end{lemma}

\begin{proof}
We bound the jump odometer \(J\) from iid Poisson input in an arbitrary finite killed square \(V\). Take a deterministic square \(K\subset V\) or the intersection of such a square with \(V\). Write \(g_x=G_K(\cdot,x)\), continued by zero outside \(K\), and define \(W(x)=\E_x[J(X_{\tau_K})]\), with \(J=0\) off \(V\). The expectation is over an auxiliary walk started at \(x\) and killed on leaving \(K\), with the boundary values of the realized field \(J\) held fixed. Mass balance for the jumps, with initial counts \(\sigma\) and final counts \(\eta\), gives
\begin{equation}\label{e:mass}
 J(x)=W(x)+\langle g_x,\sigma-\eta\rangle+\mathcal M_x .
\end{equation}
Indeed, the changes of \(g_x\) along the jumps telescope to \(\langle g_x,\eta-\sigma\rangle\). Their conditional means at origins in \(K\) sum to \(-J(x)\). At an exterior origin \(y\), the mean is \(\tfrac14\sum_{u\in K:\,u\sim y}g_x(u)=\Prob_x(X_{\tau_K}=y)\), by Green symmetry and counting the final step of the killed walk. Thus the exterior means sum to \(W(x)\). The remaining term \(\mathcal M_x\) is the sum of centered jump increments. For the whole square \(V\), \(W=0\).

We record uniform estimates for a test square of radius at most \(r\). The mean exit time is \(O(r^2)\), since the lazy normal coordinate in an interval of length \(b\) has mean time \(2i(b+1-i)\) from position \(i\). The killed Green diagonal is \(O(\log(2+r))\). To see this, use \(G_K=\int_0^\infty e^{-t(I-Q_K)}\,dt\). The free coordinate return bound \(C(1+t)^{-1/2}\) follows by integrating \(e^{-t(1-\cos\theta)/2}\), so the two-dimensional return integral up to \(r^2\) is \(O(\log(2+r))\). At time \(r^2\) each kernel entry is \(O(r^{-2})\); after that the survival probability decays exponentially on scale \(r^2\), by iterating the mean-exit bound and Markov's inequality. Clipping the square can only reduce these upper bounds. By the first-hit formula, \(\|g_x\|_\infty=g_x(x)\), and by symmetry \(\sum g_x\le Cr^2\). The killed energy is
\[
 \calE(g,g)=\langle g,(I-Q_K)g\rangle
       =\tfrac14\sum_{\{y,z\}}(g(y)-g(z))^2 ,
\]
where we count every unoriented interior edge and each exterior edge once, with zero exterior value. In particular \(\calE(g_x,g_x)=g_x(x)\).

Let \(M\) be dyadic and \(r=\floor{a_0\sqrt M}\) for a small fixed \(a_0>0\). Reveal the fact of a jump before its direction and keep only the first \(3M\) jumps from each origin in the martingale. On the cap \(J\le3M\) throughout the test square and its neighboring sites, this is the actual martingale. Its conditional variance sum is at most \(CM\log(2+r)\), by the energy bound, and its increments are bounded by \(C\log(2+r)\). If a martingale has conditional variance sum at most \(A\) and increments bounded by \(b\), the probability of a deviation \(h\) in either direction is at most \(2e^{-c h^2/(A+bh)}\). This follows by Taylor expansion of the conditional exponential at \(|s|b\le1\), iteration, and optimization. Deterministic cutoffs and Fatou transfer the bound to the stopped sum. A deviation of size \(\epsilon r^2\) therefore fails with probability at most \(C e^{-c_\epsilon M/\log(2+M)}\) on the cap, also for squares of radius a fixed positive fraction of \(r\). The weighted Poisson input has the same estimate for its difference from \(\rho\sum g_x\), from its generating function and \(\sum g_x^2\le\|g_x\|_\infty\sum g_x\). These estimates may be unioned over \(O(M)\) deterministic poles.

Call \(z\) an approximate maximum at scale \(M\) if \(j=J(z)\in[M,2M)\) and every site within distance \(r+1\) in the maximum norm has \(J\le(1+\delta_0)j\). We bound the probability of this event uniformly in \(V,z\). Use \(K=V\cap(z+[-r,r]^2)\). By \eqref{e:mass}, on the good deviations, \(|J-W|\le C_\rho r^2\) throughout \(K\), with \(\delta_0\le1/4\) so the jump truncation agrees. If \(z\) is within \(c_0r\) of an exterior face of \(V\), there is a fixed positive probability \(\theta_0\) of exiting \(V\) before leaving the \(r\)-window. For example, in \(\floor{\epsilon_0r^2}\) steps the normal coordinate has variance of order \(\epsilon_0r^2\) and fourth moment of order \(\epsilon_0^2 r^4\), so symmetry and Cauchy--Schwarz give a fixed lower probability of a displacement past the face once \(c_0\) is sufficiently small relative to \(\sqrt{\epsilon_0}\). The stopped second-moment maximal bound makes the chance of leaving the \(r\)-window by this time \(O(\epsilon_0)\). Choose \(\epsilon_0\) first. The exit contribution is then at most \((1-\theta_0)(1+\delta_0)j\), contradicting \(J(z)=j\) for small \(a_0,\delta_0\).

Otherwise \(W\) is nonnegative and harmonic on the concentric open square of radius \(c_0r/2\). It satisfies \(W(z)\ge j-C_\rho r^2\). The discrete Harnack inequality \cite[Theorem~6.3.9]{LawlerLimic} gives \(W\ge c_1j\) on the closed concentric square of radius \(c_0r/4\). The simple walk satisfies its hypotheses and the ratios of the two square sizes are fixed. Thus \(J\ge c_1j/2\) there on the good deviations after reducing \(a_0\). Take \(K_0=z+[-\floor{c_0r/16},\floor{c_0r/16}]^2\) and tile it up to strips of width \(O(\ell)\) by odd squares of side \(\ell\asymp M^{3/8}\). The source of any port into these tiles lies within one step of \(K_0\), inside the region \(J\ge c_1j/2\) of radius \(c_0r/4\) for large \(r\). The deterministic prefix of \(\floor{c_1M/4}\) jump directions at each such source includes \(c_2M\) jumps of each direction with exponentially small failure for a suitable \(c_2>0\), by the binomial exponential bound. In that inner square the actual inlet thus contains \(c_2M\) particles per port and has total input per port at most \(3M\).

For any deterministic inlet in this range choose \(B=\ceil{C_2\ell^2}\) extra particles per port. There are enough since \(\ell^2=o(M)\). Use them for an internal diffusion-limited aggregation prefix. Move each selected label until it parks at a site where no earlier label parked, or exits. At a site with a parked label a moving trial is never alone, so its sleep attempts are ineffective. This is a legal prefix for the whole local input.

To estimate cover failure, couple each trial to a full prospective killed walk continued after parking. Fix a site \(y\). Let \(A_y\) count the full walks that hit \(y\) and \(B_y\) count those whose auxiliary post-parking continuations hit \(y\). Each count uses at most one hit per label. If \(y\) remains unfilled then \(A_y\le B_y\). The independent paths for \(A_y\) give mean \(4B/G(y,y)\), since the number \(d(x)\) of outward ports obeys \((I-Q)\1=d/4\) and the hit probability is \(G(x,y)/G(y,y)\). Conditional on the stopped histories, \(B_y\) is dominated by the number of hits from one independent killed walk per vertex, whose mean is \(O(\ell^2)/G(y,y)\). A threshold between these two means, with \(C_2\) large, gives two one-sided Bernoulli-sum bounds. Union over \(y\) to bound cover failure by \(C\ell^2 e^{-c\ell^2/\log(2+\ell)}\).

Only the stopped legal prefix determines the cover event. On success, at least one active remains per site and the auxiliary post-parking paths have not been used in the actual stabilization. By the saturation-plus-input identity in Section~\ref{s:finite}, the continuation on the unread stacks has terminal law \(\pi_\ell\) conditionally on this stopped legal history.

For each tile the stationary tail assumption and the cover bound can now be unioned over its at most \((3M+1)^{4\ell}\) deterministic inlet rows. Both probability exponents dominate \(O(\ell\log M)\). Union also over the tiles before identifying any actual inlet. The restriction identity then gives at least \(v\ell^2\) final sleepers in each such tile, outside a stretched-exponentially rare event at scale \(M\).

Use \(g_z=G_{K_0}(\cdot,z)\) in \eqref{e:mass}. Since \(0\le\eta\le1\), replacing \(g_z\) by its average on a tile costs at most the sum of absolute deviations on that tile. Summed over the tiles, the cost is \(O(r\ell\sqrt{\log(2+r)})\) after summing the absolute errors. In fact the tile squared error is at most \(C\ell^2\) times its nearest-neighbor energy, by two applications of one-dimensional Poincaré, and the summed \(L^1\) error is bounded by Cauchy--Schwarz and the energy of \(g_z\). The discarded strips contribute \(O(r\ell\log(2+r))\). Both are \(o(r^2)\). The sum of \(g_z\) is at least \(c_3r^2\), by the second-moment maximal bound up to a small multiple of \(r^2\) steps from the center of \(K_0\). Thus the weighted final occupation is at least \(v\sum g_z-o(r^2)\). The weighted Poisson input and centered jump errors can have tolerance less than a fixed small multiple of \((v-\rho)r^2\), by the deviation estimates. It follows that
\[
 j\le (1+\delta_0)j-c_4(v-\rho)r^2.
\]
Since \(r^2\asymp a_0^2M\), choose \(\delta_0>0\) small enough to contradict this inequality. We have proved that an approximate maximum at a deterministic center and scale has probability at most \(C e^{-c M^\beta}\), with \(\beta>0\), uniformly over the ambient size.

Finally start at the origin if \(J(0)\ge M_0\), and at each center that is not an approximate maximum move to a point within its window with \(J\) larger by factor \(1+\delta_0\). In a finite stabilization this procedure ends. An endpoint of dyadic scale \(M\) is within \(C_{\delta_0}\sqrt M\) of the origin, by summing the previous move lengths along the backwards geometric progression. Union over \(O_{\delta_0}(M)\) deterministic possible endpoint sites and all dyadic \(M\ge M_0/2\). This bounds the tail of \(J(0)\) by a quantity tending to zero as \(M_0\to\infty\), uniformly over \(V\).

In the infinite Poisson field use one draw of the field independent of the whole array, restricting both to each killed box. A smaller box's stabilization is legal in a larger box if its exterior is held frozen. Least action gives monotone exhaustion odometers. The origin jump odometer has finite limit by the uniform bound. Infinitely many entries of the origin stack are jumps almost surely, so finiteness of the limiting jump count implies finiteness of the limiting instruction count. The resulting product law is the one in Rolla--Sidoravicius \cite[Lemma~4]{RS}, with generating jump kernel, unit jump rate, positive sleep rate and iid finite-mean initial counts. This identifies the odometer event with almost-sure continuous-time local fixation, hence with the finite-jump event above. Their zero--one correspondence and monotonicity under Poisson thinning also make the almost-sure nonfixation set upward closed. If \(u_->\zeta\), take \(\zeta<\rho<v<u_-\). The fixed negative field for \(v\) supplies the lemma's tail with \(\alpha=1\), a contradiction.
\end{proof}

\section{Pressure at fields of order \texorpdfstring{\(L^{-2}\)}{L to the power -2}}\label{s:small}

\begin{proposition}\label{p:small}
 For every \(\kappa>0\) there is a constant \(C_\kappa\) such that for all sufficiently large odd \(L\) and \(0<a\le1/L\),
 \begin{equation}\label{e:small}
 \Phi(-a)\le\Phi_L(-a)+\kappa a+C_\kappa/L^2.
 \end{equation}
 For each \(\varepsilon>0\), let \(\mathcal L_L(\varepsilon)\) be the event
 that each translated odd side-\(s\) square in \(V_L\) with
 \(s\ge C_{\lambda,\varepsilon}\log L\) has at least
 \((\zeta-\varepsilon)s^2\) sleepers. Then for a suitable
 \(C_{\lambda,\varepsilon}\),
 \begin{align}
 \liminf_{L\to\infty}\pi_L(N/L^2<\zeta+\varepsilon)
     &\ge p_\varepsilon>0, \label{e:low}\\
 \Prob\big(\mathcal L_L(\varepsilon)\big)&\longrightarrow 1 .\label{e:local}
 \end{align}
 The probability in \eqref{e:local} is for the saturated sample.
\end{proposition}

\begin{proof}
We first control the inlet sizes and their variation for the nine tiles of side \(S\) in a saturated side-\(3S\) square. A nonanticipating choice of active label permits its next unused jump direction to be read from an independent prospective killed walk for that label. The resulting jump and sleep decisions have the stack law. The sum of the visits of all the full prospective walks bounds the actual jump count at each site. For a given target site the visits from one walk are zero or, after a hit, geometric of mean \(G(x,x)\le C\log(2S)\). At a small exponential parameter of order \(1/\log(2S)\), its exponential moment minus one is bounded by a constant times the parameter times its mean, by summing the geometric series. The sum of these means is the exit potential \(O(S^2)\). Independence of the prospective walks then gives, for \(0\le a\le1\), the cap
\[
 \max_x J(x)\le M=C_1 S^2(1+a\log S)
\]
with failure probability at most \(e^{-20aS^2-cS}\), after allowing room in the constants for the following bounds.

For bounded \(f\) on the side-\(3S\) square, a nearest-neighbor difference of \(Gf\) is bounded by \(CS\|f\|_\infty\). To verify this, integrate the killed product heat kernel for rate-\(1/2\) coordinates. For adjacent starting vertices only one coordinate row changes. The row sum of the unchanged killed factor is at most one. The odd-image formula bounds the absolute row sum of the changing coordinate by twice the absolute discrete gradient sum for the free coordinate. This last sum is twice the return kernel, by symmetry and monotonicity away from zero. Monotonicity follows by the heat equation for consecutive differences on the half-line. The bound is \(C(1+t)^{-1/2}\). For \(t>S^2\) use the bound at \(S^2\) and the exponential survival estimate for the remaining time. The integral is \(O(S)\). At a boundary edge the neighboring value is zero and the one-coordinate exit time from the first row is \(O(S)\). In particular the variation of \(G(1-\eta)\) on any line of \(S\) incident sites is \(O(S^2)\).

For the side vector \(I=(I_i)_{i=1}^S\), write \(\TV(I)=\sum_{i=1}^{S-1}|I_{i+1}-I_i|\). To bound the fluctuation in \(J=G(1-\eta)+\mathcal M\) from \eqref{e:mass}, fix signs on the \(S-1\) edges of a line and put \(c_s=\sum_{(x,x')}s_{xx'}(\delta_{x'}-\delta_x)\). The signed martingale sum of consecutive differences uses \(G c_s\). Its energy is \(\langle Gc_s,c_s\rangle\le C\sqrt S\,\calE(Gc_s,Gc_s)^{1/2}\), hence \(O(S)\). Keeping at most the first \(M\) jumps from each origin bounds the sum of conditional variances by \(CMS\) and an increment bound \(C\sqrt S\). At deviation \(C_2 S\sqrt{M(1+aS)}\), the martingale bound has exponent at least \(c C_2 S(1+aS)\) for large \(C_2\), because \(S(1+aS)/M\) is bounded. Take the union over the \(2^{S-1}\) sign patterns. The directed count into a port differs from \(J/4\) by at most the centered maximum over the first \(M\) directions at its origin site. Along one incoming side these maxima \(Z_i\) depend on distinct site lists. They satisfy \(\E e^{Z_i^2/(C M)}\le C\) by the bounded-increment exponential maximum bound. Use \((\sum_i Z_i)^2\le S\sum_i Z_i^2\) and their independence to bound the sum by \(C_3 S\sqrt{M(1+aS)}\) with the same required exponential error. The total variation of these directed errors is at most twice that sum. Thus, writing \(A_{a,S}=\sqrt{(1+a\log S)(1+aS)}\), all actual incoming side vectors satisfy
\begin{equation}\label{e:variation}
 0\le I_i\le M,\qquad \TV(I)\le CS^2 A_{a,S},
\end{equation}
except on an event of probability at most \(e^{-20aS^2-cS}\), uniformly for the field range. Exterior sides have zero inlet.

Next compare two fixed deterministic inlets \(I'\ge I\) in a side-\(S\) tile. Write \(N_I\) for the count after stabilizing the saturated input plus \(I\). Release the extra labels after this stop on the same stacks. Commuting the extra additions to the beginning identifies the output with \(N_{I'}\) on that same stack array. The additional labels move on tails fresh at the first stop. If \(1\le d\le S/4\), a new label starting on a boundary site must reach depth \(d\) before hitting the outside coordinate face to end in the depth-\(d\) core. Use the normal coordinate with value \(0\) on the exterior face and \(1\) at the entering boundary site. A free normal-coordinate walk started there hits \(d\) before \(0\) with probability \(1/d\), by the stopped linear martingale. The prospective walks of the new labels are independent conditionally on the first stopped history. Thus their potential core entrants are bounded by an independent Bernoulli sum \(Z\) of mean at most \(\Delta/d\), where \(\Delta=\sum(I'-I)\). There are at most \(CSd\) collar sites and at most \(N_I\) old labels. It follows that
\begin{equation}\label{e:top}
 N_{I'}-N_I\le CSd+Z .
\end{equation}
In particular if \(g\ge4\), \(S/g\ge2\), \(\Delta\le C_*S^3/g^2\) and \(d=\floor{S/g}\), a bound on \(\E[e^Z\mid\text{first history}]\) gives
\begin{equation}\label{e:weighted}
 \E[e^{-aN_I};\,N_{I'}-N_I>C_4S^2/g]
       \le e^{-cS^2/g}\,\E_{\pi_S}e^{-aN}.
\end{equation}
Here \(C_*,C_4\) are fixed and \(C_4\) is sufficiently large compared to \(C_*\). The first-history conditional tail bound is multiplied by the measurable first count weight. This relative estimate will pay for a union over the unrounded profiles.

For each side in \eqref{e:variation} choose a fixed partition into \(r=\min(S,\ceil{g^2 A_{a,S}})\) groups of lengths at most \(\ceil{S/r}\). Replace all entries of a group by its maximum rounded up in integer increments of \(b=\ceil{S^2/g^2}\). The increment summed on a side is at most \(Sb+\ceil{S/r}\TV(I)\) if \(r<S\), and at most \(Sb\) if \(r=S\). Thus the increment over each tile is \(O(S^3/g^2)\). The total log size of the rounded joint table is bounded by
\[
 C(1+g^2 A_{a,S})\log\!\big(2+Cg^2(1+a\log S)\big).
\]
Here the groups are fixed before seeing the inlets and each group height has at most \(Cg^2(1+a\log S)\) possible rounded values.

Write \(Z_S(a)=\E_{\pi_S}e^{-aN}\) and divide all contributions by \(Z_S(a)^9\). The work and variation exception contributes at most \(e^{-11aS^2-cS}\) because \(Z_S(a)\ge e^{-aS^2}\).

On the cap, sum the event of a bad top-up comparison over the unrounded rows satisfying \eqref{e:variation}. Their number is at most \((M+1)^{36S}\). Fix one such row and drop the event that it is the actual row. For each choice of a bad tile apply \eqref{e:weighted} there, and multiply the eight independent tile factors \(Z_S(a)\) from the rest of that row. The resulting relative contribution is at most \(9e^{CS\log S-cS^2/g}\).

On the remaining event \(N_I\ge N_{I'}-C_4S^2/g\) on every tile. Since the field is negative, \(e^{-aN_I}\le e^{aC_4S^2/g}e^{-aN_{I'}}\). This estimate can be summed over the smaller set of rounded rows. Drop each selection event before taking the product of the independent rounded tile moments. Taking logarithms gives, for \(S/g\) larger than a fixed multiple of \(\log S\),
\begin{equation}\label{e:triple}
 \begin{aligned}
 \Phi_{3S}(-a)&\le\Phi_S(-a)+\frac{C a}{g}+\frac{Ce^{-c'S}}{S^2}\\
 &\quad+\frac{C(1+g^2 A_{a,S})\log(2+Cg^2(1+a\log S))}{S^2}.
 \end{aligned}
\end{equation}

Fix a field \(0<a\le1/L\) and take \(S_i=3^iL,\ g_i=g_0(i+1)^2\). For all sufficiently large \(L\) the condition on \(S_i/g_i\) holds uniformly in \(i\). Choose \(g_0\) with \(\sum C/g_i\le\kappa\). Since \(aS_i\le3^i\) and \(1+a\log S_i\le C(i+1)\), the entropy series is bounded by \(C_\kappa L^{-2}\): each term is at most \(L^{-2}\) times a polynomial in \(i+1\) times \(3^{-3i/2}\). The exception series has the same bound. The limit through these odd larger sizes at the same fixed field is \(\Phi(-a)\), which proves \eqref{e:small}.

Convexity and Lemma~\ref{l:fix} give \(\Phi(-a)\ge-a\zeta\). Now choose \(a=D/L^2\), with \(D\) fixed, after choosing \(\kappa<\varepsilon/4\). For \(D\) large enough, \eqref{e:small} yields
\[
 \E_{\pi_L} e^{-DN/L^2}\ge e^{-D(\zeta+\kappa)-C_\kappa}
       >2e^{-D(\zeta+\varepsilon)} .
\]
The contribution on \(N/L^2\ge\zeta+\varepsilon\) is at most \(e^{-D(\zeta+\varepsilon)}\), proving \eqref{e:low}. In the other direction \eqref{e:lower} gives the upper limit \(e^{-D\zeta}\) for this expectation at any fixed \(D>0\). Dividing \eqref{e:small} by \(-D/L^2\) gives
\[
 \frac{\Phi(-D/L^2)}{-D/L^2}
   \ge -D^{-1}\log \E_{\pi_L} e^{-DN/L^2}-\kappa-C_\kappa/D.
\]
Send \(L\to\infty\) to obtain \(u_-\ge\zeta-\kappa-C_\kappa/D\). Send \(D\to\infty\), then \(\kappa\downarrow0\), so \(u_-=\zeta\).

For \(v<\zeta\) a fixed negative field now bounds \(\pi_s(N<v s^2)\le e^{-c_v s^2}\) for all sufficiently large odd \(s\). Within the actual saturation of \(V_L\) impose a polynomial workload cap, with failure tending to zero by the exit potential estimate. Each translated side-\(s\) square has at most \(\exp(Cs\log L)\) possible inlet profiles under this cap. A union over the deterministic profiles and the at most \(O(L^3)\) squares proves \eqref{e:local} by taking \(v=\zeta-\varepsilon\) and \(s\ge C_{\lambda,\varepsilon}\log L\). When \(v\le0\) the claim is automatic.
\end{proof}

\section{Green profiles of the low-density samples}\label{s:profiles}

For saturation of \(V_L\) let \(J_L\) be the jump count field, \(\operatorname{Dir}_L(v,e)\) the number of jumps from \(v\) in direction \(e\), and \(\eta_L\) the final occupations. Put \(h=(L+1)^{-1}\), \(\Omega=(0,1)^2\), and set
\[
 u_L(hv)=h^2G_{V_L}(1-\eta_L)(v).
\]
Set \(u_L=0\) at the boundary mesh points and interpolate bilinearly. With probability tending to one the following inequalities hold with deterministic errors tending to zero:
\begin{equation}\label{e:field}
 \begin{gathered}
 \max J_L\le B L^2,\qquad \max_v |h^2J_L(v)-u_L(hv)|=o(1),\\
 \max_{v,e}|h^2\operatorname{Dir}_L(v,e)-u_L(hv)/4|=o(1).
 \end{gathered}
\end{equation}
The full walk visit estimate in Proposition~\ref{p:small} at \(a=0\), applied on \(V_L\) itself, gives the cap. For the second bound truncate at the cap in the whole-square Green martingale \eqref{e:mass}. Each pole gives variance bound \(CL^2\log(2L)\) and increments bounded by \(C\log(2L)\). A deviation of \(L^{3/2}\) has probability at most \(C e^{-c L/\log(2L)}\). For the last bound the maximal fluctuation of each of the first \(BL^2\) entries in a direction list at this deviation has probability at most \(2e^{-cL}\). Union over the poles and directions, and multiply by \(h^2\).

The row sum bound \(O(L^2)\) and adjacent row-difference bound \(O(L)\) from Proposition~\ref{p:small} apply to \(G_{V_L}(1-\eta_L)\). Multiplication by \(h^2\) bounds the interpolated fields uniformly and makes each mesh difference \(O(h)\). This gives the Lipschitz bound up to the boundary. Take any sequence of configurations satisfying \eqref{e:field}, \eqref{e:local}, and \(N/L^2<\zeta+\varepsilon\). Along a subsequence the empirical measures \(h^2\sum_v\eta_L(v)\delta_{hv}\) on the closed square converge weakly to \(\varrho\,dx\), \(0\le\varrho\le1\), since they are dominated by the grid measures converging to Lebesgue measure. Any fixed interior square may be approximated by translated odd grid squares, to which \eqref{e:local} applies. Differentiation over squares gives \(\varrho\ge\zeta-\varepsilon\) almost everywhere. As the total mass is at most \(\zeta+\varepsilon\),
\[
 \|\varrho-\zeta\|_1\le 3\varepsilon,\qquad
 \|\varrho-\zeta\|_2^2\le3\varepsilon .
\]
Take a further subsequence with uniform limit \(u_L\to u\). Test the discrete Green equation against a smooth compactly supported function. By symmetry and \(h^{-2}(I-Q)\varphi\to-\Delta\varphi/4\), we obtain the zero Dirichlet equation \(-\Delta u/4=1-\varrho\). Write \(\calG=(-\Delta_\Omega/4)^{-1}\) with zero boundary, \(u_*=(1-\zeta)\calG1\), and \(w=u_*-u=\calG(\varrho-\zeta)\). The sine basis with eigenvalues \(\pi^2(m^2+n^2)/4\) gives
\begin{equation}\label{e:hess}
 \|D^2w\|_2\le C\sqrt\varepsilon.
\end{equation}
Indeed a pure or mixed second derivative introduces a multiplier bounded by \(C(m^2+n^2)\), and the corresponding sine and cosine systems are orthogonal. The undifferentiated expansion is uniformly convergent by Cauchy--Schwarz and summability of \((m^2+n^2)^{-2}\). Its continuous zero-boundary representative agrees with the uniform limit, by uniqueness of the distributionally harmonic difference with zero boundary values. Also \(\|w\|_\infty\le C\).

We construct the approximating set at two resolutions. Fix a tolerance \(0<\tau<1\) and odd integers \(l,\gamma\ge3\). Set \(H=1/l,\ k=\gamma l,\ q=\tau/(10k^2)\). Round each value of \(w\) on the coarse mesh of spacing \(H\) to the nearest multiple of \(q\), and call the bilinear interpolant \(W\). On a coarse cell \(R\) the distance in supremum norm to a suitable affine function is at most \(CH\|D^2 w\|_{2,R}\). To verify the scaling, subtract an affine function with the same mean and mean gradient, apply Poincaré to the gradient and to the function on the unit square, and use the supremum bound from integrating the fundamental theorem of calculus in both coordinates. The bound applies by approximation to \(W^{2,2}\) functions and also to a rectangle of fixed aspect ratio. Bilinear interpolation reproduces the affine part with convex weights. The error after rounding is at most \(C H\|D^2w\|_{2,R}+q\). Thus cells with error more than \(\tau/k^2\) have fraction bounded by
\begin{equation}\label{e:bad}
 C\gamma^4\|D^2w\|_2^2/\tau^2.
\end{equation}

To count the arrays, record the first two integers of each row and then its horizontal second differences. At the unrounded level each second difference annihilates affine functions and is bounded by \(CH\|D^2w\|_{2,P}\) for a patch \(P\) of width \(2H\) and height \(H\), using an interior patch along boundary rows. These patches have bounded overlap. Cauchy--Schwarz over \(O(l^2)\) positions gives total absolute unrounded second difference at most \(C\|D^2w\|_2\). The encoded rounded differences therefore have total absolute sum at most \(C l^2+Cq^{-1}\|D^2w\|_2\). If \(\gamma^2\|D^2w\|_2\le\tau\), this is \(O(l^2)\). The number of possible integer differences is at most \(e^{C_0 l^2}\) by counting magnitudes of given total and then signs. The first two entries of each row have \(O(q^{-1})\) choices, since \(\|w\|_\infty\le C\). Keep arrays with this difference budget and the uniform bound on every node. This gives a deterministic catalogue \(\mathcal W_k\) with
\begin{equation}\label{e:entropy}
 \log|\mathcal W_k|\le C_0l^2+C_0l\log(2k/\tau),
\end{equation}
where \(C_0\) is independent of \(\tau,l,\gamma\).

Fix \(\tau,\gamma\) and a fraction \(\theta>0\). Choose \(\varepsilon>0\) small enough that \eqref{e:hess} implies the difference-budget condition and the right side of \eqref{e:bad} is at most \(\theta\). This choice is independent of \(l\). For each fixed \(k=\gamma l\), partition a lattice square of side \(L=kS\) into the corresponding \(k^2\) side-\(S\) tiles. For sufficiently large odd \(S\), every sample on the low-density event with local bound \eqref{e:local} and jump-field bounds \eqref{e:field} admits a \(W\in\mathcal W_k\) such that all but at most \(\theta k^2\) tiles satisfy the following bound for the actual count \(I\) of arrivals along every incoming boundary edge with receiving vertex \(v\):
\begin{equation}\label{e:fit}
 \left|(L+1)^{-2} I-\tfrac14(u_*-W)\big((v-(1/2,1/2))/L\big)\right|
      \le \tau/k^2 .
\end{equation}
For proof, any contrary sequence at fixed \(k\) has the uniform field limit just constructed. Its rounded array fits the limit on the good coarse cells to error \(\tau/k^2\), and each bad coarse cell excludes only \(\gamma^2\) tiles. The indicated ideal coordinate lies in the receiving fine tile and differs from the emitting site's \(h\)-mesh coordinate by \(O(1/L)\). The directed count bound in \eqref{e:field} and uniform convergence then make the left side of \eqref{e:fit} at most \(\tau/(4k^2)+o(1)\) on each good tile, with \(k\) fixed. Outer ports have zero inlet and require no approximation.

\section{The upper density bound}\label{s:upper}

Suppose there are \(0<\delta<1/8\) and an odd sequence \(S\to\infty\) with \(S^{-2}\E_{\pi_S}N\ge\zeta+4\delta\). We compare the saturated side-\(kS\) sample to stabilizations with fixed arrival profiles in its tiles. Independence across tiles will be used for each fixed family of profiles. For fixed \(W\), define the integer upper inlet vector \(I^+(W,j)\) on tile \(j\). At an interior-facing port with receiving vertex \(v\), put \(\xi=(v-(1/2,1/2))/L\) and use the entry
\[
 \max\!\left\{0,\,
 \ceil{(L+1)^2\big((u_*-W)(\xi)/4+2\tau/k^2\big)}\right\}.
\]
Assign zero at outer ports. By \eqref{e:fit}, its excess above the actual inlet in every good tile is nonnegative and bounded by \(3\tau(L+1)^2/k^2+1\) per port. For large \(S\) this is at most \(C\tau S^2\).

For any fixed pair \(I,I^+\) with this much excess per port, \(\Delta\le C\tau S^3\). Recall from \eqref{e:top} that adding the excess arrivals costs at most the \(CSd\) collar sites and a Bernoulli sum of mean at most \(\Delta/d\). At \(d=\floor{\sqrt\tau S}\) both bounds are \(O(\sqrt\tau S^2)\). Choose \(\tau\) small enough, depending on \(\delta\), that \(\Prob(N_{I^+}-N_I>\delta S^2)\le e^{-c_\delta S^2}\) by the Bernoulli exponential bound. The jump cap from \eqref{e:field} gives at most \((B L^2+1)^{4S}\) deterministic inlet rows per tile. At fixed \(k\), union over these rows, all \(k^2\) tiles, and the finite catalogue, considering only pairs within the stated excess. The failure probability tends to zero. This statement is simultaneous over the deterministic rows before the actual inlet and the fitted \(W\) are selected.

Choose odd \(\gamma\) large with \(C_0/\gamma^2<\delta^2/2\), set \(\theta=\delta\), and choose \(\varepsilon<\delta/2\) small enough for the catalogue with these fixed \(\tau,\gamma\). Equation \eqref{e:low} fixes \(p_\varepsilon>0\) before choosing \(l\). Take an odd \(l\) sufficiently large that \(k=\gamma l\) satisfies \(\log|\mathcal W_k|\le\delta^2 k^2\) and \(e^{-\delta^2 k^2}<p_\varepsilon/2\). From now on \(k\) is fixed while \(S\) tends to infinity along the supposed high-mean sequence.

Write \(N_L=\sum_v\eta_L(v)\) for the ambient sleeper count. The ambient run and all the finitely many deterministic tests complete almost surely at each size. On the low event \(N_L/L^2<\zeta+\varepsilon\) and the simultaneous good events, take the fitted \(W\) and its good tiles. By restriction the true tile count is the count for its exact inlet on its own stacks. On each good tile it is at least \(N_{I^+(W,j)}-\delta S^2\), and the missing test count on a bad tile is at most \(S^2\). Thus
\begin{equation}\label{e:sum}
 \frac{N_L}{L^2}\ge
    \frac1{k^2}\sum_{j=1}^{k^2}\frac{N_{I^+(W,j)}}{S^2}
       -\theta-\delta .
\end{equation}
For each fixed \(W\) the test counts have independent laws from \(\pi_S\) on disjoint stacks. They are evaluated without conditioning on the event that \(W\) fits the ambient sample. Their densities lie in \([0,1]\) with means at least \(\zeta+4\delta\). The threshold \(\zeta+\varepsilon+\theta+\delta\) is at least \(\delta\) below these means. The bounded-variable exponential estimate gives probability at most \(e^{-2\delta^2 k^2}\) that their average is below this threshold. Indeed the centered log moment has second derivative equal to a tilted variance at most \(1/4\). Union over the deterministic catalogue and use \eqref{e:sum} to bound the low-event probability by \(e^{-\delta^2k^2}+o(1)\), where the error tends to zero with \(S\) at this fixed \(k\). This contradicts \eqref{e:low} at \(L=kS\). No such high-mean sequence exists. Together with \eqref{e:lower} this proves Theorem~\ref{t:main}.

\section*{Acknowledgment}
OpenAI language models assisted with exploring the argument, checking the proof and preparing the manuscript. The author takes responsibility for the mathematical claims and their presentation.

\begin{thebibliography}{9}
\bibitem{Forien} N. Forien, \emph{A new proof of superadditivity and of the density conjecture for Activated Random Walks on the line}, 2025, \href{https://arxiv.org/abs/2502.02579v2}{arXiv:2502.02579v2}.
\bibitem{HJJ} C. Hoffman, T. Johnson and M. Junge, \emph{The density conjecture for activated random walk}, 2024, \href{https://arxiv.org/abs/2406.01731}{arXiv:2406.01731}.
\bibitem{LawlerLimic} G. F. Lawler and V. Limic, \emph{Random Walk: A Modern Introduction}, Cambridge Studies in Advanced Mathematics, vol.~123, Cambridge University Press, 2010.
\bibitem{LL} L. Levine and F. Liang, \emph{Exact sampling and fast mixing of Activated Random Walk}, \href{https://arxiv.org/abs/2110.14008v3}{arXiv:2110.14008v3}, 2024.
\bibitem{RS} L. T. Rolla and V. Sidoravicius, \emph{Absorbing-state phase transition for driven-dissipative stochastic dynamics on \(\mathbb Z\)}, \href{https://arxiv.org/abs/0908.1152v2}{arXiv:0908.1152v2}, 2011.
\bibitem{RT} L. T. Rolla and L. Tournier, \emph{Non-fixation for biased activated random walks}, Ann. Inst. Henri Poincar\'e Probab. Stat. \textbf{54} (2018), 938--951, \href{https://doi.org/10.1214/17-AIHP827}{doi:10.1214/17-AIHP827}.
\end{thebibliography}
\end{document}
